\documentclass{beamer}

\usetheme{Antibes}
\usepackage[utf8]{inputenc}
\usepackage{amsmath}
\usepackage{amssymb}
\usepackage{amsfonts}
%%%%%%%%%%%%%%%%%%%%%%%%%%%%%%%%%%
%%%% New commands below
%%%% math fonts
\usepackage{etoolbox,xparse,dsfont,bm,amssymb}
\newcommand{\myMFabc}[4]{\expandafter#1\csname#3#4\endcsname{{#2{#4}}}}
%\forcsvlist{\myMF{\newcommand}{\bm}{bm}}{x,y}
\newcommand{\myMFcmd}[4]{\expandafter#1\csname#3#4\endcsname{{#2{\csname#4\endcsname}}}}
%\forcsvlist{\myMFcmd{\newcommand}{\bm}{bm}}{theta,phi}


\newcommand{\MFabc}[3][\newcommand]{
    \def\doOld##1##2{\forcsvlist{\myMFabc{#1}{##1}{##2}}{#3}}
    \providecommand{\do}{do}
    \RenewDocumentCommand \do { >{\SplitList{,}} m } { \doOld##1 }
    \docsvlist{#2}
}
% \MFabc{ {\bm,bm}, {\mathcal,cal} }{A,B,C}
\newcommand{\MFcmd}[3][\newcommand]{
    \def\doOld##1##2{\forcsvlist{\myMFcmd{#1}{##1}{##2}}{#3}}
    \providecommand{\do}{do}
    \RenewDocumentCommand \do { >{\SplitList{,}} m } { \doOld##1 }
    \docsvlist{#2}
}

%%% 0 and 1
\newcommand{\bmzero}{{\bm{0}}}
\newcommand{\bmone}{{\bm{1}}}
\def\bbone{{\ensuremath{\mathds{1}}}}
\let\one\bbone
%  mathfrak, e.g., \frakT = \mathfrak{T}
\MFabc{ {\mathfrak,frak}, }{T,u,U,o,O,E,m}
%  mathbb, e.g., \R = \mathbb{R}
\MFabc{ {\mathbb, }, }{A,N,Z,R,C,Q}

%%% all capital letters
% \mathcal, \mathscr, and \mathbb 
% e.g., \calA=\mathcal{A}
% \bm\mathcal, e.g., \calbmA = \bm{\mathcal{A}}
\newcommand{\hatbm}[1]{\widehat{\bm{#1}}}
\newcommand{\tildebm}[1]{\widetilde{\bm{#1}}}
\newcommand{\bmcal}[1]{\bm{\mathcal{#1}}}
\newcommand{\caltilde}[1]{\mathcal{\widetilde{#1}}}
\newcommand{\calhat}[1]{\mathcal{\widehat{#1}}} 
\newcommand{\scrtilde}[1]{\widetilde{\mathscr{#1}\mspace{1mu}\mspace{-1mu}}}% \mspace{1mu}\mspace{-1mu} to avoid spacing conflict \mathscr and \widetilde
\newcommand{\scrhat}[1]{\mathscr{\widehat{#1}\mspace{1mu}\mspace{-1mu}}}
\newcommand{\bmcalhat}[1]{\bm{\mathcal{\widehat{#1}}}}
\newcommand{\bmcaltilde}[1]{\bm{\mathcal{\widetilde{#1}}}}
%%%%%%
\MFabc{ {\mathcal,cal}, {\mathscr,scr}, {\mathbb,bb}, {\bmcal,bmcal}, {\bmcal,calbm}, {\caltilde,caltilde}, {\caltilde,tildecal}, {\calhat,calhat}, {\calhat,hatcal}, {\scrtilde,scrtilde}, {\scrtilde,tildescr}, {\scrhat,scrhat}, {\scrhat,hatscr}, {\bmcalhat,bmcalhat}, {\bmcalhat,bmhatcal}, 
{\bmcalhat,calbmhat}, {\bmcalhat,calhatbm}, {\bmcalhat,hatbmcal}, {\bmcalhat,hatcalbm}, {\bmcaltilde,bmcaltilde}, {\bmcaltilde,bmtildecal}, {\bmcaltilde,calbmtilde}, {\bmcaltilde,caltildebm}, {\bmcaltilde,tildebmcal}, {\bmcaltilde,tildecalbm}}{A,B,C,D,E,F,G,H,I,J,K,L,M,N,O,P,Q,R,S,T,U,V,W,X,Y,Z}

\renewcommand{\scrtildeA}{\widetilde{\mathscr{A}\mspace{2mu}}\mspace{-6.1mu}}% \mspace{2mu} and\mspace{-6.1mu} to avoid spacing conflict \mathscr and \widetilde
\renewcommand{\scrhatA}{\mathscr{\widehat{A}\mspace{2mu}}\mspace{-6.1mu}}
\let\tildescrA\scrtildeA
\let\hatscrA\scrhatA

%%% all letters
% bm, \widehat， \widetilde, and
% \mathsf shortcuts -- for random variables， 
% e.g., \bmA=\bm{A}, \bmalpha=\bm{\alpha}
%%%%%%%%%
\MFabc{{\bm,bm}, {\mathsf,sf}, {\widehat,hat}, {\widetilde,tilde}, {\hatbm,hatbm}, {\hatbm,bmhat}, {\tildebm,tildebm}, {\tildebm,bmtilde}}{a,b,c,d,e,f,g,h,i,j,k,l,m,n,o,p,q,r,s,t,u,v,w,x,y,z,A,B,C,D,E,F,G,H,I,J,K,L,M,N,O,P,Q,R,S,T,U,V,W,X,Y,Z}

%%% all Greek alphabet, e.g., \bmalpha = \bm{\alpha}
\let\eps\varepsilon
\MFcmd{{\bm,bm}, {\widehat,hat}, {\widetilde,tilde}, {\tildebm, tildebm}, {\tildebm, bmtilde}, {\hatbm, hatbm}, {\hatbm, bmhat}}{alpha,beta,gamma,delta,epsilon,zeta,eta,theta,iota,kappa,lambda,mu,nu,xi,omicron,pi,rho,sigma,tau,upsilon,phi,chi,psi,omega,Alpha,Beta,Gamma,Delta,Epsilon,Zeta,Eta,Theta,Iota,Kappa,Lambda,Mu,Nu,Xi,Omicron,Pi,Rho,Sigma,Tau,Upsilon,Phi,Chi,Psi,Omega,varrho,varphi,vartheta,varepsilon,varsigma,ell,eps}



%%%%%%%%%% activation functions
%\newcommand{\actdef}[1]{\lowercase{\expandafter\def\csname#1\endcsname}{\texttt{#1}}}
\newcommand{\actdef}[1]{\expandafter\def\csname#1\endcsname{{\ensuremath{\mathtt{#1}}}}}
\forcsvlist{\actdef}{ReLU, LReLU, LeakyReLU, ELU, GELU, SiLU, Softplus, dGELU, dSiLU, dSoftplus, Tanh, Sigmoid, Arctan, Softsign, SRS, dSRS, Swish, dSwish, Mish, dMish, SELU, CELU, dSELU, Sin,SinLU, SinTU, PSinTU, sine, cosine, Sine, Cosine, EUAF}

%%%%%%%%%%%%%%%%%%%%%%%

\title{Neural Tangent Kernel Analysis of Multi-component and Multi-layer Neural Networks (MMNNs)}
\author{Based on the provided report}
\date{\today}

\begin{document}

\begin{frame}
\titlepage
\end{frame}

\begin{frame}{Introduction \& Goal}
\begin{itemize}
    \item \textbf{Objective:} To understand the local optimization landscape of MMNNs.
    \item \textbf{Method:} Analyze the Neural Tangent Kernel (NTK) in the infinite-width limit.
    \item \textbf{Core Hypothesis (S1):} Only the external linear layer parameters ($\tilde{\bm{A}}$) are trained. The internal parameters ($\tilde{\bm{W}}$) defining the basis functions are randomly initialized and fixed.
\end{itemize}
\end{frame}

\begin{frame}{MMFN Architecture: Single Layer}
A single layer is a linear combination of random basis functions:
\begin{equation*}
    \bmh(\bm{x}) = \bm{A} \sigma(\bm{W}\bm{x} + \bm{b}) + \bm{c}
\end{equation*}
\vspace{1em}
Using an augmented representation to absorb the biases:
\begin{equation*}
    \bmh(\bm{x}) = \tilde{\bm{A}}
    \begin{bmatrix}
        \sigma(\tilde{\bm{W}}\tilde{\bm{x}}) \\
        1
    \end{bmatrix}
    \quad \text{where} \quad
    \begin{cases}
      \tilde{\bm{x}} = [\bm{x}^{\top}, 1]^{\top} \\
      \tilde{\bm{W}} = [\bm{W}, \bm{b}] \\
      \tilde{\bm{A}} = [\bm{A}, \bm{c}]
    \end{cases}
\end{equation*}
\begin{itemize}
    \item Each output $\bmh[i]$ is a \textit{component}.
    \item The fixed $\sigma(\tilde{\bm{W}}\tilde{\bm{x}})$ acts as a basis.
    \item The trainable $\tilde{\bm{A}}$ learns the coefficients for that basis.
\end{itemize}
\end{frame}

\begin{frame}{Inspiration: Function Decomposition}
\begin{itemize}
    \item MMNNs are inspired by a "divide and conquer" strategy for complex functions.
    \item \textbf{N-Dimensional Decomposition Theorem:} Any function on a hyperrectangle can be decomposed via the principle of inclusion-exclusion. The 1D case is the simplest version of this.
    \begin{equation*}
    \resizebox{\hsize}{!}{$f(\bm{x}) = \sum_{S \subseteq \{1, \dots, d\}} (-1)^{d-|S|} \sum_{\bm{j}_{\bar{S}}} \left( \mathcal{D}_S \left[ f(\cdot, \bm{x}_{\bar{S}, \bm{j}_{\bar{S}}}) \right] (\bm{x}_S) \right)$}
    \end{equation*}
    \item The fixed weights $\tilde{\bm{W}}$ create a basis analogous to partitioning the domain (the operators $\mathcal{D}_S$).
    \item The trainable weights $\tilde{\bm{A}}$ learn the coefficients, analogous to approximating the simpler function pieces on the sub-domains.
\end{itemize}
\end{frame}

\begin{frame}{Analysis Setup: Key Assumptions}
\begin{block}{Activation Function}
    We focus on \textbf{ReLU}, as its positive-homogeneity simplifies kernel computations.
\end{block}

\begin{block}{Initialization Strategy ($n \to \infty$)}
\begin{itemize}
    \item \textbf{Inner Weights $\tilde{\bm{W}}$:} Drawn from an order-one distribution. This is crucial for creating a stable partition of the input space.
    \item \textbf{Outer Weights $\tilde{\bm{A}}$:} Scaled by $1/\sqrt{n}$. This is essential for the NTK theory, ensuring convergence to a well-behaved Gaussian Process.
\end{itemize}
\end{block}
\end{frame}

\begin{frame}{Single-Layer NTK: From Sum to Expectation}
\begin{alertblock}{Focus of this Analysis}
The following derivation is exclusively for a \textbf{single-layer MMFN}. The multi-layer case will be discussed at the end as future work.
\end{alertblock}
The kernel is the sum of gradient products over trainable parameters ($\{A_j\}, c$). For a single output MMFN:
\begin{equation*}
    h(\bm{x}) = \frac{1}{\sqrt{n}} \sum_{j=1}^n A_j \sigma(\bm{w}_j^\top \bm{x} + b_j) + c
\end{equation*}
In the infinite-width limit ($n \to \infty$), the Law of Large Numbers gives:
\begin{equation}
    K_{NTK}(\bm{x}, \bm{x}') = \mathbb{E}_{\bm{w},b} \left[ \sigma(\bm{w}^\top \bm{x} + b)\sigma(\bm{w}^\top \bm{x}' + b) \right] + 1
\end{equation}
\end{frame}

\begin{frame}{Integral Formulation \& Parameter Mapping}
The expectation is an integral over the parameter space of $(\bm{w}, b)$.
\begin{itemize}
    \item We perform a change of variables from the high-dimensional parameter space to the 2D pre-activation space $(g_1, g_2)$.
    \item This maps the problem to an expectation over a 2D bivariate Gaussian.
\end{itemize}
\begin{block}{Bivariate Gaussian Parameters}
\begin{itemize}
    \item \textbf{Means:} $\mu = \mu_b$
    \item \textbf{Variances:} $\sigma_1^2 = \bm{x}^\top \Sigma_w \bm{x} + \sigma_b^2$
    \item \textbf{Covariance:} $\text{Cov}(g_1, g_2) = \bm{x}^\top \Sigma_w \bm{x}' + \sigma_b^2$
\end{itemize}
\end{block}
\end{frame}

\begin{frame}{Comparison with Standard FCNN Kernel}
\begin{itemize}
    \item The NTK for a standard 1-layer FCNN often assumes centered pre-activations: $\mathbb{E}_{\bm{w}}[\sigma(\bm{w}^\top\bm{x})\sigma(\bm{w}^\top\bm{x}')]$.
    \item \textbf{The key difference for MMFN:} The random, non-trainable bias $b$ is \textit{inside} the activation function.
    \item This makes the pre-activations non-centered (if $\mu_b \neq 0$) and leads to a more complex, but richer, kernel calculation.
\end{itemize}
\end{frame}

\begin{frame}{The Final Result: Analytical NTK}
\begin{theorem}[NTK for a Single-Layer ReLU MMFN]
The expectation $\mathbb{E}[\text{ReLU}(g_1)\text{ReLU}(g_2)]$ for jointly Gaussian pre-activations $(g_1, g_2)$ is:
\begin{equation*}
\resizebox{0.9\hsize}{!}{$\mathbb{E}[\sigma(g_1)\sigma(g_2)] = \sigma_1\sigma_2 \cdot I_{12} + \mu_1\sigma_2 \cdot I_{01} + \mu_2\sigma_1 \cdot I_{10} + \mu_1\mu_2 \cdot I_{00}$}
\end{equation*}
The terms $I_{ij}$ are moments of standardized bivariate normal variables $(Z_1, Z_2)$ with thresholds $a_{1,2} = -\mu_{1,2}/\sigma_{1,2}$ and correlation $\rho$.
For the symmetric case ($a_1=a_2=a$), these terms are:
\begin{itemize}
    \item $I_0 = \mathbb{E}[\mathbf{1}_{Z>a, Z'>a}] = L(-a, -a; \rho)$
    \item $I_1 = \mathbb{E}[Z \cdot \mathbf{1}_{Z>a, Z'>a}] = (1+\rho)\phi(a)\Phi\left(a\sqrt{\frac{1-\rho}{1+\rho}}\right)$
    \item $I_2 = \mathbb{E}[Z Z' \cdot \mathbf{1}_{Z>a, Z'>a}] = \rho L(-a, -a; -\rho) + \phi(a, a; \rho)$
\end{itemize}

\end{theorem}
\end{frame}

\begin{frame}{Final Notes}
\vspace{0.5em}
where $\phi(z) = \frac{1}{\sqrt{2\pi}}e^{-z^2/2}$, $\Phi(z) = \int_{-\infty}^z \phi(t)dt$, and $L(a,b;\rho) = \int_{-\infty}^a\int_{-\infty}^b \phi(z_1,z_2;\rho)dz_1dz_2$.
\end{frame}

\begin{frame}{Future Work: The Multi-Layer Case}
\begin{itemize}
    \item The analysis for a deep MMNN builds upon the single-layer kernel via a recursive formula.
    \item The kernel at layer $L$ depends on the kernel from layer $L-1$.
    \item The detailed computation for this recursive case, especially for ResMMNNs, is significantly more complex and is left for future work.
    \item This single-layer derivation serves as the essential base case.
\end{itemize}
\end{frame}

\end{document}
